In this section we point out how to generalized \cref{cor:SSofFFCat} to a case where we replace the framing structure on the flow category $\fC$ by other tangential structures.
As a result we get a spectral sequence that calculates the bordism homology $\Omega_*^{\cX}(\norm{\fC})$.

To start with, we will generalize the definition of normal framing.
\begin{defn}
    A tangential structure is a collection of fibrations $\cX(n) \to BO(n)$, with maps $\cX(n) \to \cX(n+1)$ such that
    \begin{equation*}
        \xymatrix
        {
            \cX(n) \ar[r]\ar[d]  &  \cX(n+1) \ar[d] \\
            BO(n) \ar[r]                    & BO(n+1)
        }
    \end{equation*}
\end{defn}

\begin{defn}
    Let $M \hookrightarrow \Hlf^k \times \bbR^N$ be an embedded $m$ dimensional $\k$ manifold. A \textbf{normal $\cX$ structure on $M$} is a choice of a lift
    \begin{equation}
        \xymatrix
        {
              &&  \const{\cX(N)} \ar[d] \\
            M \ar[rr]_{\hspace{-10pt} \nu M} \ar@{-->}[urr]^{\Tilde{\nu M}}  && \const{BO(N+k-m)}
        }
    \end{equation}
    where the diagram takes place in the category $\k$ spaces, which is a functorial $\op^k$ family of lifts
    \begin{equation}
        \xymatrix
        {
              &&  \cX(N) \ar[d] \\
            M(A) \ar[rr]_{\hspace{-10pt}\nu M(A)} \ar@{-->}[urr]  && BO(N+k-m)
        }
    \end{equation}
    for each $A \subseteq \op^k$.
\end{defn}
\begin{defn}
    We say that an $\cX$ structure is \textbf{multiplicative} if it is equipped with structure maps
    \begin{equation*}
        \cX(n) \times \cX(m) \to \cX(n+m)
    \end{equation*}
    that commutes with the maps
    \begin{equation*}
        BO(n) \times BO(m) \to BO(n+m)
    \end{equation*}
\end{defn}
Note that if $\cX$ is multiplicative then a choice of $\cX$ structure on two embedded manifolds $M,L$ induces a $\cX$ structure on $M\times L$ in the following way
\begin{equation*}
\begin{tikzcd}
	& {\const{\cX(d) \times \cX(d')}} & {\const{\cX(d+d')}} \\
	{L \times M} & {\const{BO(d)\times BO(d')}} & {\const{BO(d+d')}}
	\arrow[from=1-2, to=1-3]
	\arrow[from=1-2, to=2-2]
	\arrow[from=1-3, to=2-3]
	\arrow[dashed, from=2-1, to=1-2]
	\arrow[from=2-1, to=2-2]
	\arrow[from=2-2, to=2-3]
\end{tikzcd}
\end{equation*}
where $d$ and $d'$ are the codimensions of $L$ and $M$, respectively.

For an tangential structure $\cX$ we denote by $\zeta_n$ the universal real vector bundle over $X(n)$, which is (by abuse of notation) the pullback of the universal real vector bundle $\zeta_n \to BO(n)$.
Let $M$ be an $\cX$ structure manifold (for simplicity - a closed manifold) and denote by $d$ the codimension of $M$ in the ambient space, we obtain the following diagram:
\begin{equation*}
    \begin{tikzcd}
	{\nu M} & {\zeta_n} & {\zeta_n} \\
	M & {X(d)} & {BO(d)}
	\arrow[from=1-1, to=1-2]
	\arrow[from=1-1, to=2-1]
	\arrow["\lrcorner"{anchor=center, pos=0.125}, draw=none, from=1-1, to=2-2]
	\arrow[from=1-2, to=1-3]
	\arrow[from=1-2, to=2-2]
	\arrow["\lrcorner"{anchor=center, pos=0.125}, draw=none, from=1-2, to=2-3]
	\arrow[from=1-3, to=2-3]
	\arrow["\Tilde{\nu M}",from=2-1, to=2-2]
	\arrow[from=2-2, to=2-3]
\end{tikzcd}
\end{equation*}
where the right square is pullback by definition and the same for the whole rectangle. We denote
\begin{equation}\label{eq:StrAnalogOfFraming}
    \varphi_M \colon \nu M \to \zeta_n
\end{equation}
Note that although we omitted $\Tilde{\nu M}$ from the notation, a different choices of $\Tilde{\nu M}$ defines different map $\varphi_M$.
Note that if $M$ is $\k$ manifold then $\varphi_M$ is a map of $\k$ vector bundles.

Given a multiplicative structure we can define an $\cX$ structured flow category (compare to \cref{def:NormalFramingOfFlowCat}).
\begin{defn}
    Let $\fC$ be an embedded flow category, an $\cX$ structure on $\fC$ is a choice of $\cX$ structure on each embedded morphism manifold $M_{y,x}$ such that 
    \begin{equation*}
         \partial_y \varphi_{z,x}= \varphi_{z,y} \oplus \varphi_{y,x}
    \end{equation*}
    where $\varphi_{y,x}$ is the isomorphism of $\<{J_{y,x}}$ bundles arise from the $\cX$ structure on $M_{y,x}$, see \cref{eq:StrAnalogOfFraming}.
\end{defn}

Similarly to \cref{con:topChOfCat}, using Pontryagin Thom we get a coherent chain complex in $\Mod_{M\cX}$ where $M\cX$ is the Thom spectrum of $\cX$.
\begin{con}
    Let $\fC$ be a bounded $\cX$ flow category.
    For each $x,y \in \fC$ we define a map
    \begin{equation*}
        g_{y,x}: S^{N(\norm{y}-\norm{x})}\smashSpaces (\Hlf^{J_{y,x}})_+ \to  Th(\nu M_{y,x}) \to Th(\zeta_{N(\norm{y}-\norm{x})}).
    \end{equation*}
    where the first map is the collapse map of the embedding $e_{y,x}$ (see \cref{def:CollapaseMap}) and the second map is $Th(\varphi_{y,x})$.
    The coherence of being $\cX$ flow categories amounts to
    \begin{equation*}
        g_{z,y} \smashSpaces g_{y,x} = \partial_y g_{z,x}
    \end{equation*}
    as maps between diagrams in $\Top$.

    By taking $\Sigma^{N(\norm{y}-\norm{x})} \Sinf g_{y,x}$ we get a collection of maps to the Thom spectrum of $\cX$
    \begin{equation*}
        \bbS \otimes (\Hlf^{J_{y,x}})_+ \to  M\cX.
    \end{equation*}
    which is equivalent to a map
    \begin{equation*}
        (\Hlf^{J_{y,x}})_+ \to \Hom_{\Sp^O}(\bbS, M\cX),
    \end{equation*}\todo{I am confuse, we can move between strict functors of topologically enriched categories to $\infty$ functors in $\Sp$, but can we do the same for $M \cX$ modules?}
    the last step is to use the adjunction between $\bbS-\Mod$ and $M\cX-\Mod$ to obtain a map
    \begin{equation*}
        \bar{g}_{y,x}:(\Hlf^{J_{y,x}})_+ \to \Hom_{M \cX}(M\cX, M\cX).
    \end{equation*}\todo{HEre is a question: Let $E \in \Sp^O$ do we have a category of $E$ modules and adjoinction?}
    To get that the collection of these maps define a functor
    \begin{equation*}
        \J \to \Mod-\cX
    \end{equation*}
    we use the following commutative diagram:
    \begin{equation*}
        % https://q.uiver.app/#q=WzAsNixbMCwwLCJcXEhvbV97XFxUb3B9KFNeZCxcXHpldGFfZClcXHRpbWVzIFxcSG9tX3tcXFRvcH0oU157ZCd9LFxcemV0YV97ZCd9KSJdLFswLDEsIlxcSG9tX3tcXFNwfShcXGJiUyxNIFxcY1gpIFxcdGltZXMgXFxIb21fe1xcU3B9KFxcYmJTLE0gXFxjWCkiXSxbMCwyLCJcXEhvbV97TSBcXGNYfShNIFxcY1gsTSBcXGNYKSBcXHRpbWVzIFxcSG9tX3tNIFxcY1h9KE0gXFxjWCxNIFxcY1gpIl0sWzIsMiwiXFxIb21fe00gXFxjWH0oTSBcXGNYLE0gXFxjWCkiXSxbMiwxLCJcXEhvbV97XFxTcH0oXFxiYlMsTSBcXGNYKSJdLFsyLDAsIlxcSG9tX3tcXFRvcH0oU157ZCtkJ30sXFx6ZXRhX3tkK2QnfSkiXSxbMCw1LCJcXHdlZGdlIl0sWzUsNCwiKFxcU157LShkK2QnKX1cXFNpbmYiXSxbMCwxLCJcXFNeey1kfVxcU2luZiBcXFNeey1kJ31cXFNpbmYiLDJdLFsyLDMsIlxcY2lyYyIsMl0sWzQsMywiXFxNb2QiLDJdLFsxLDIsIlxcTW9kIiwyXV0=
        \begin{tikzcd}
	{\Hom_{\Top}(S^d,\zeta_d)\times \Hom_{\Top}(S^{d'},\zeta_{d'})} && {\Hom_{\Top}(S^{d+d'},\zeta_{d+d'})} \\
	{\Hom_{\Sp}(\bbS,M \cX) \times \Hom_{\Sp}(\bbS,M \cX)} && {\Hom_{\Sp}(\bbS,M \cX)} \\
	{\Hom_{M \cX}(M \cX,M \cX) \times \Hom_{M \cX}(M \cX,M \cX)} && {\Hom_{M \cX}(M \cX,M \cX)}
	\arrow["\wedge", from=1-1, to=1-3]
	\arrow["{(\S^{-d}\Sinf ,\S^{-d'}\Sinf)}"', from=1-1, to=2-1]
	\arrow["{\S^{-(d+d')}\Sinf}", from=1-3, to=2-3]
	\arrow["\Mod"', from=2-1, to=3-1]
	\arrow["\Mod", from=2-3, to=3-3]
	\arrow["\circ"', from=3-1, to=3-3]
        \end{tikzcd}
    \end{equation*}
    where $\Mod$ denote the $\bbS-\Mod, \cX-\Mod$ adjunction. We obtain that
    \begin{equation*}
        \bar{g}_{z,y} \smashSp \bar{g}_{y,x} = \partial_y \bar{g}_{z,x}
    \end{equation*}
    hence we got a (topologically enriched) coherent chain complex $\Fun(\J,\Mod-\cX)$.
\end{con}

The definition of a right flow module (\cref{def:RightFlowModule}) is modified by replacing framed flow modules by $\cX$ structured flow modules and $\Omega^{\fr}$ by $\Omega^{\cX}$.


Similarly, we define when two $\cX$ flow categories are composable. To define the composition category we modify \cref{def:CompCat} we get a boundary diagram
\begin{equation*}
    D \colon \op^{\{b,\dots a\}} \setminus \const{0} \to \Top
\end{equation*}
such that $D(\{a\})$ are $\cX$ manifolds with $\<{A\setminus\{a\}}$ corner structure.
he lifts of the structure map $\nu D(B)$ assemble into a lift of the structure map $\nu \colim D$ hence induces a $\cX$ structure on the colimit. We get that
\begin{cor}
    Let $\fC, \sD$ be two composable flow categories bounded by $b+1, a-1$, with $\cX$ structure and flow manifolds written as $M_{-,-},N_{-,-}$ then
    \begin{equation*}
          \bigcup_{b \leq p \leq a} M_{b+1,p} \times N_{p,a-1} \coloneqq \colim D
    \end{equation*}
    is an $\cX$ manifold (compare to \cref{def:CompCat}).
    Where $D$ is defined in the following way: for each $B = \{i_k,\dots i_1\} \subseteq \{b,\dots a\}$
    \begin{equation*}
        D(i_k,\dots i_1) = M_{b+1,i_k}M_{i_k,i_{k-1}} \times \dots M_{i_2,i_1} \times N_{i_1,a-1}.
    \end{equation*}
    We write
    \begin{equation*}
        M_{b+1,a-1} \coloneqq \bigcup_{b \leq p \leq a} M_{b+1,p} \times N_{p,a-1} \coloneqq \colim D.
    \end{equation*}
\end{cor}
We define composition category\todo{resolve (what is the comp. cat. ?)} as same as the framed case.

\ThA*
\begin{rem}
    The convergence of the spectral sequence in \cref{SSofStrFFCat} relies on the convergence results established for chain complexes (\cref{SSofCh}). Specifically, if the flow category $\fC$ is bounded below, it induces a bounded below chain complex in $\Ch(\Mod_{\cX})$, which by \cref{cor:strong_cond_conv} yields strong convergence. If we omit the bounded below condition on $\fC$, we still obtain an exact couple that conditionally converges to the generalized homology of the realization, as detailed in \cref{cor:strong_cond_conv}.
\end{rem}

\subsection{Change of rings}
Give a framed flow category $\fC$, we can benefit from the fact that the changing the tangential structure of manifolds correspond to maps between spectra, to get a spectral sequence that calculates the bordism homology (represented by Thom spectra) of a manifold.

Let $\cX,\cY$ be two tangential structures.
Denote by $\FlowX$ (resp. $\FlowY)$ the collections of $\cX$ (resp. $\cY$ ) structured flow categories.

A map $\cX \to \cY$ induce a "forgetting" map $F$, from $\cX$ manifolds to $\cY$ manifolds. The fact that Cartesian product commutes with $F$, induces a map:
\begin{equation*}
    (-)^{\cY}: \FlowX \to \FlowY
\end{equation*}
On the other hand, the map $\cX \to \cY$ induces a map of spectra $M\cX \to M\cY$ and hence an extension of scalars of coherent chain complexes:
\begin{equation*}
    \otimes_{M\cX} M\cY \colon \Ch(\Mod_{M\cX}) \to \Ch(\Mod_{M\cY})
\end{equation*}
induces by post composition
\begin{equation*}
    \Ch \xrightarrow[]{K} \Mod_{M\cX} \xrightarrow[]{\otimes_{M\cX} M\cY} \Mod_{M\cY}
\end{equation*}
\begin{lem}\todo{Re-write the proof}
        The following diagram commutes:
        \begin{equation*}
            % https://q.uiver.app/#q=WzAsNCxbMCwwLCJcXEZsb3dYIl0sWzAsMSwiXFxGbG93WSJdLFsxLDEsIlxcQ2goTW9kX3tNXFxjWX0pIl0sWzEsMCwiXFxDaChNb2Rfe01cXGNYfSkiXSxbMCwxXSxbMSwyXSxbMCwzXSxbMywyXV0=
            \begin{tikzcd}
            	\FlowX & {\Ch(Mod_{M\cX})} \\
            	\FlowY & {\Ch(Mod_{M\cY})}
            	\arrow["{K_{(-)}}"', from=1-1, to=1-2]
            	\arrow["(-)^{\cY}"', from=1-1, to=2-1]
            	\arrow["{\otimes_{M\cX} M\cY}", from=1-2, to=2-2]
            	\arrow["{K_{(-)}}"', from=2-1, to=2-2]
            \end{tikzcd}
        \end{equation*}
\end{lem}
\begin{proof}
        Let $\fC$ be a $\cX$ flow category. On the one hand we have $K_{\fC^{\cY}}$, on the other $K_{\fC}\otimes_{M\cX} M\cY$.
        The object of the coherent chains are:
        \begin{equation*}
            (K_{\fC^{\cY}})_p =
            \WedgeSp_{\norm{z}=p} M\cY = 
            \WedgeSp_{\norm{z}=p} M\cX \otimes_{M\cX} M\cY \simeq 
            (K_{\fC} \otimes_{M\cX} M\cY)_p.
        \end{equation*}
        We have that the maps of $K_{\fC^{\cY}}$ are the collapse maps:
        \begin{equation}\label{eq:Y_chainMap}
            c_{M_{z,w}^{\cY}}: \Sinf \Hlf^{J_{z,w}} \to M\cY
        \end{equation}
        where  $M^{\cY}$ is the manifold $M$ with $\cY$ structure.
        The maps of $K_{\fC}\otimes_{M\cX} M\cY$ are given by collapsing and tensoring
        \begin{equation}\label{eq:Xchanged_chainMap}
            c_{M_{z,w}}\otimes_{M\cX} M\cY: \Sinf \Hlf^{J_{z,w}} \to M\cX \xrightarrow{(-)\otimes_{M\cX} M\cY} M\cY.
        \end{equation}
        The fact that the $\cY$ structure of $M_{z,w}$ is coming from $\cX$ yields that the collapse map of $M_{z,w}^{\cY}$ factors through $\cX$, the functoriality of the collapse map give the equivalence between \cref{eq:Y_chainMap} and \cref{eq:Xchanged_chainMap}.
\end{proof}%\todo{can be rewrited more clearly. For example just claim that $c_N \otimes_{M\cX} M\cY \simeq c_{N^{\cY}}$ for any manifold $N$}

First, we get that we can start with framed flow category and forget to other tangential structure, to get that:
\begin{cor}
    Let $\fC$ be a framed flow category and let $\cX$ be a tangential structure. Then the spectral sequence of \cref{SSofStrFFCat} implies to $\fC^{\cX}$ having:
    \begin{equation*}
        E_r \implies 
        \pi_*(\mid \fC^{\cX} \mid) \simeq
        \pi_*(\norm{\fC} \otimes M\cX)
        \eqqcolon M\cX_*(\norm{\fC})
    \end{equation*}
\end{cor}
The first equivalence holds because the realization from $\Ch$ to $\Fil$ is given by colimit, hence commutes with tensoring (as left adjoint).

Second, we can apply it to Morse function, to calculate the bordism homology of a manifold:
\begin{cor}\label{cor:SSofMorseStr}
    Let $f:M\to \bbR$ be a Smale Morse function and let $\cX$ be a tangential structure, then the spectral sequence of the framed flow category $(\fC_{f})^{\cX}$ converge to $M\cX_*(M)$.
\end{cor}
\begin{rem}
    Since the Morse index is inherently bounded below (by $0$), the induced flow category $\fC_{f}$ is bounded below. Thus, \cref{cor:strong_cond_conv} guarantees strong convergence to the bordism homology.
\end{rem}